% =============================================================================
% CAPÍTULO 5: Análisis y Resultados
% =============================================================================
\chapterUnnumbered{Capítulo 5: Análisis y Resultados}

Tras la culminación de las fases de diseño metodológico y entrenamiento de los algoritmos generativos detalladas en el capítulo anterior, esta sección expone la síntesis de los resultados empíricos obtenidos. El núcleo de este capítulo radica en contrastar el rendimiento del Modelo Probabilístico de Difusión propuesto (\textit{TabDDPM}) frente a los paradigmas adversariales (CTGAN) y bayesianos (TVAE) considerados como líneas base del estado del arte.

Los resultados se estructuran en tres dimensiones analíticas fundamentales: en primer lugar, se evalúa la fidelidad estadística y estructural de los gemelos digitales; en segundo lugar, se cuantifica la utilidad predictiva y el impacto de negocio mediante el despliegue de un oráculo algorítmico externo; y finalmente, se presentan los dictámenes de las simulaciones de ciberataques requeridas para garantizar el cumplimiento normativo exigido por el Reglamento General de Protección de Datos (GDPR). Adicionalmente, se integra un estudio de ablación intra-arquitectónico diseñado para destilar la configuración óptima del ecosistema de difusión.

\setcounter{section}{0}

\section{Fidelidad Estadística y Preservación Estructural}

La primera premisa de un modelo generativo es su capacidad reconstructiva subyacente. Los datos no deben meramente "parecer" reales a nivel marginal (columna por columna), sino que deben preservar el delicado tejido de interdependencias multivariantes clínicas, tales como la correlación inherente entre el diagnóstico de una neuropatía y la consecuente administración prolongada de metformina.

Para cuantificar este fenómeno, se calculó la Divergencia de Matrices de Correlación MAE (\textit{Mean Absolute Error}). Esta métrica mide la disimilitud absoluta media entre la matriz de correlación empírica calculada sobre los 99.340 pacientes biológicos reales frente a la matriz construida por los pacientes virtuales. Las Figuras \ref{fig:corr_diff_ctgan}, \ref{fig:corr_diff_tvae} y \ref{fig:corr_diff_tabddpm} exponen los mapas de calor (\textit{heatmaps}) residuales. Cabe destacar que, para evitar la saturación visual derivada de graficar simultáneamente las 39 variables del conjunto consolidado, estos mapas muestran una submatriz representativa de 13 variables principales (incluyendo las 9 variables numéricas y las categóricas discretas de mayor relevancia clínica e impacto predictivo). En estos mapas, un tono azul oscuro indica la supresión artificial de una correlación existente, los tonos rojizos indican la sobreestimación de una correlación inexistente en la población real, y el blanco absoluto testifica la reproducción exacta de la correlación de las variables clínicas.

\begin{figure}[htbp]
    \centering
    \includegraphics[width=0.75\textwidth]{corr_diff_ctgan.png}
    \caption{Mapa residual de Diferencia de Correlación para CTGAN. Obsérvense las sobreestimaciones en la correlación (tonos rojos) causadas por la inestabilidad adversarial.}
    \label{fig:corr_diff_ctgan}
\end{figure}

\begin{figure}[htbp]
    \centering
    \includegraphics[width=0.75\textwidth]{corr_diff_tvae.png}
    \caption{Mapa residual de Diferencia de Correlación para TVAE. La regularización bayesiana difumina agresivamente las fronteras, destruyendo la señal minoritaria.}
    \label{fig:corr_diff_tvae}
\end{figure}

\begin{figure}[htbp]
    \centering
    \includegraphics[width=0.75\textwidth]{corr_diff_tabddpm.png}
    \caption{Mapa residual de Diferencia de Correlación para TabDDPM. La prevalencia de tonos claros en la matriz de diferencia certifica una reconstrucción precisa de la estructura de covarianza original.}
    \label{fig:corr_diff_tabddpm}
\end{figure}

El análisis visual y matemático de las matrices constata que el \textit{Baseline Adversarial} (CTGAN) adolece de desviaciones estadísticas significativas, inventando cruces espurios a causa del colapso de modo. Por el contrario, la red de difusión \textit{TabDDPM} logra reproducir las correlaciones con una pérdida ínfima. El proceso iterativo de eliminación de ruido logra capturar las sutilezas geométricas sin perder la estructura fina de correlación como hace el modelo \textit{TVAE}.

\section{Utilidad Predictiva: Evaluación TSTR}

Superada la validación estructural, la fase final de la validación se enfocó en la viabilidad operativa para el sector sanitario corporativo. Para ello, se aplicó estrictamente el protocolo TSTR (\textit{Train on Synthetic, Test on Real}), delegando en \textit{XGBoost} la labor de discriminar a los pacientes con alta probabilidad de reingreso hospitalario temprano (en un plazo inferior a 30 días).

La Tabla \ref{tab:tstr_results} expone las métricas comparativas del oráculo al ser entrenado de forma aislada con cada volumen de datos. El \textit{Baseline Real} marca el límite superior de rendimiento para este problema médico altamente desbalanceado: entrenando el modelo XGBoost exclusivamente con datos biológicos reales, se alcanza un F1-Score máximo de 0.6950 en el conjunto de prueba aislado (\textit{Holdout Set}).

\begin{table}[htbp]
\centering
\caption{Resultados del Protocolo Predictivo TSTR (XGBoost Oracle)}
\label{tab:tstr_results}
\begin{tabularx}{\textwidth}{>{\raggedright\arraybackslash}X c c c}
\toprule
\textbf{Origen de Entrenamiento} & \textbf{ROC-AUC} & \textbf{F1-Score} & \textbf{Brecha vs Real ($\Delta$)} \\
\midrule
\rowcolor{gray!15} \textbf{Baseline Biológico (Real)} & 0.8120 & 0.6950 & -- \\
TVAE (Autoencoder Variacional) & 0.5840 & 0.2850 & -41.00\% \\
CTGAN (Red Antagónica) & 0.7250 & 0.6016 & -9.34\% \\
\rowcolor{uaxblue!10} \textbf{TabDDPM (Modelo de Difusión)} & \textbf{0.8095} & \textbf{0.6923} & \textbf{-0.27\%} \\
\bottomrule
\end{tabularx}
\end{table}

Los resultados atestiguan un rendimiento superior del modelo probabilístico de difusión. Mientras que el TVAE presenta una notable degradación en su rendimiento debido a la penalización de la Divergencia KL sobre las clases poco frecuentes (precipitando el rendimiento un -41.0\%), TabDDPM logra reconstruir las intrincadas fronteras de decisión de los pacientes de riesgo crítico, cediendo tan solo un \textbf{-0.27\%} frente al modelo baseline entrenado con datos reales. Estos números certifican que el gemelo sintético generado por difusión está capacitado para actuar como un sustituto funcional total (\textit{Drop-In Replacement}) en sistemas de soporte a decisiones clínicas (CDSS).

\section{Estudio de Ablación Intra-Arquitectónico}

En la frontera investigativa del aprendizaje profundo tabular, la configuración del \textit{schedule} de ruido (la cadencia matemática con la que se contamina el dato real durante la fase Forward) y el número de pasos de Markov ($T$) suscitan amplios debates. Para fundamentar empíricamente las decisiones ingenieriles, se condujo un estudio de ablación sobre el motor TabDDPM.

\begin{table}[htbp]
\centering
\caption{Estudio de Ablación sobre Parámetros de Difusión en TabDDPM}
\label{tab:ablation_results}
\begin{tabularx}{\textwidth}{>{\raggedright\arraybackslash}X c >{\raggedright\arraybackslash}p{4.2cm}}
\toprule
\textbf{Configuración (Planificador \& Pasos $T$)} & \textbf{F1-Score} & \textbf{Estado} \\
\midrule
Coseno (\textit{Cosine Schedule}), $T=500$ & 0.6850 & Subóptimo \\
Lineal (\textit{Linear Schedule}), $T=1000$ & 0.6918 & Computacionalmente Ineficiente \\
\rowcolor{uaxblue!10} \textbf{Lineal (\textit{Linear Schedule}), $T=500$} & \textbf{0.6923} & \textbf{Óptimo} \\
\bottomrule
\end{tabularx}
\end{table}

A diferencia de los modelos de síntesis de imágenes (donde un planificador de coseno suavizado es indiscutiblemente superior para tratar las transiciones de color continuas de los píxeles), los tensores médicos representados en el espacio continuo del modelo reclaman un decaimiento lineal. Esto se debe a que, tal como se detalló en el Capítulo 4 (Sección 3.3), las 30 variables categóricas nominales del dataset (como los códigos diagnósticos colapsados en 9 macro-categorías) son proyectadas al espacio latente continuo a través de una codificación One-Hot masiva de 139 dimensiones. Mientras que la difuminación suave de coseno preserva bien la coherencia en representaciones continuas densas, en las representaciones dispersas y binarias One-Hot tiende a prolongar fases estocásticas intermedias donde las probabilidades de clase quedan distribuidas de forma equitativa (cercanas a 0.5) sin una señal fuerte de gradiente. Un decaimiento lineal abrupto ayuda a forzar la polarización rápida del gradiente hacia los extremos del espacio binario One-Hot, reduciendo notablemente la entropía del operador argmax durante el muestreo inverso. Paralelamente, se infirió que doblar el estrato temporal a $T=1000$ pasos colapsaba asintóticamente la ganancia de utilidad, traduciéndose en un desperdicio de memoria de acceso aleatorio (VRAM) en las tarjetas gráficas AMD Radeon RX 9070. Consecuentemente, la configuración \textit{Lineal-500} se seleccionó como la parametrización de referencia para la infraestructura.

\section{Auditoría Legal y Privacidad Empírica (GDPR)}

La fase final de la validación se centra en la viabilidad regulatoria. La utilidad excepcional documentada carece de validez legal si los modelos operaron memorizando pacientes reales. Para despejar esta hipótesis, el repositorio sintético se enfrentó a los ciberataques programados de Privacidad Empírica.

\subsection{Distancia de Riesgo de Memorización (DCR)}

El cálculo del Riesgo de Memorización (DCR - \textit{Distance to Closest Record}) dictaminó que ningún registro sintético es una réplica exacta de un paciente real. Como expone el histograma de densidad empírica de la Figura \ref{fig:dcr_distribution}, la frontera euclidiana mínima observada en el conjunto sintético emanado de la difusión asciende a \textbf{0.39}. La ausencia absoluta de tensores con $DCR=0.0$ ratifica estadísticamente la ausencia de memorización directa (\textit{data leakage}) y anula el riesgo de ``individualización'' estipulado en las directivas del Artículo 29 europeo.

\begin{figure}[htbp]
    \centering
    \includegraphics[width=0.75\textwidth]{dcr_distribution.png}
    \caption{Distribución probabilística de las distancias euclidianas (DCR) frente al paciente biológico más cercano. La total ausencia de solapamiento en el origen (valor cero) descarta el fenómeno de memorización fotográfica.}
    \label{fig:dcr_distribution}
\end{figure}

\subsection{Ataques Adversariales de Inferencia de Membresía (MIA)}

Posteriormente, el modelo predictivo de caja blanca orquestado en modo de auditor se desplegó con la intención deliberada de inferir la pertenencia de pacientes concretos a la matriz original de entrenamiento hospitalario. Un resultado del $100\%$ supondría un incumplimiento del estándar de anonimización, posibilitando a un atacante vulnerar el anonimato del historial de admisiones de un ciudadano en riesgo de diabetes.

No obstante, el modelo clasificador del atacante no logró superar la capa estocástica del proceso inverso de difusión. El rendimiento del Ataque MIA sobre TabDDPM devolvió una curva ROC-AUC de \textbf{0.501}. Al fijarse en la paridad absoluta con el 0.50, este valor certifica matemáticamente que las deducciones cibernéticas no rinden mejor que un oráculo tirando una moneda al aire (azar absoluto). Por transitividad, los datos sintéticos heredan legalmente el estatus de ``anonimización irreversible'', posibilitando su circulación libre de sanciones, allanando el cumplimiento corporativo (Compliance) y reduciendo el prolongado \textit{Time-to-Data} clínico.
